%% ============================================================================
\section{Excluded Coins and Candidate Episodes}
\label{app:excl}
Three further zero-SVB-exposure coins exist: USDP, GUSD, and BUSD. The exclusion
criterion is calm-window behavior alone. Over a calm, pre-SVB baseline (February
16--March 8, 2023), USDP, GUSD, and BUSD already print below \$0.99 in 7, 1, and 1 of 480
hours. The $1\%$ break test therefore has no discriminating power for them, and we exclude them
before reading any crisis-window price. For completeness, USDP, GUSD, and BUSD trough at \$0.948, \$0.928,
and \$0.985. Liquidity offers no alternative exclusion rule. Frozen
median daily volumes range from ${\sim}\$16$M for USDP, which was listed on Binance
mid-window, to ${\sim}\$687$M for BUSD. GUSD has no freezable series. BUSD is further
confounded by an unrelated regulatory event, the February 2023 NYDFS order halting new
issuance.
Anadu et al.~\cite{anadu2024} find GUSD and USDP had no significant dollar flows. Their
finding corroborates independently that the run was coin-specific. Three placebo coins whose
feeds can register a $1\%$ break support the test better than a six-coin panel in which
three cannot.

The $1\%$ break threshold and the worst-case floor are distinct: a coin can cross the
$1\%$ materiality threshold without nearing the price at which no reserve loss could
rationalize the price. The bound is evaluated at each hourly price, so breaching hours need not be contiguous. USDC's impossibility window is interrupted at 08{:}57 by
a recovery to \$0.9216, a price 16 bps above the floor, at which $q^{\star}=0.98$, before falling back
below the floor for seven further hours. USDT's worst in-window hour reaches $q^{\star}=0.095$, an
order of magnitude below the $q^{\star}=1$ firing threshold.

\paragraph{TrueUSD and the public-information scope condition.} The scope condition of
Section~\ref{sec:disc} rests on a documented case. TrueUSD is excluded on timing: its
impairment was quantified publicly only after our window closed. The SEC's September 2024
complaint alleges that a substantial fraction of TrueUSD's backing sat in an offshore
commodity fund through our window~\cite{sectruecoin2024}. The complaint became public
eighteen months after the March 2023 episode had ended. Precondition~(i) requires public
quantification while the exposure is unresolved, and TUSD had none. TrueUSD's
contemporaneous attestations used unquantified ``cash and cash equivalents'' language and
never named the fund. Run on TUSD in March 2023, the estimator would have drawn $\phi=0$
from those attestations and correctly returned null despite a then-unquantified impairment. TUSD
therefore illustrates the scope condition and enters no test, because its de-peg depth
over 2022--2024 is uncorrelated with its alleged impairment trajectory.

\paragraph{The seven candidate episodes.} BUSD and
Tether~(2022) are genuine price dislocations with no disclosed own-reserve
impairment. Paxos affirmed full $1{:}1$ backing throughout while redeeming ${\sim}\$7.9$B
of BUSD in a month~\cite{paxos2023busd}, and Tether's disclosures give asset-composition
breakdowns and denials of rumors without stating an impairment fraction. Run on either, the
estimator sets $\phi=0$ and reproduces the degenerate Terra null
(Section~\ref{sec:sensitivity}). Neither gives a second identification.
DAI~(2020, ``Black Thursday'')~\cite{Gudgeon_2020} has a disclosed bad-debt figure but is a
collateral-liquidation cascade requiring a different fundamental-value model, violating
precondition~(iii). USDR is a tokenized-real-estate hybrid violating~(A1). USDP and GUSD
fail the diagnostic-print test above. TUSD fails precondition~(i).

